% E3 methodology snippet, to be inserted after the E1/E2 observational methods.
\paragraph{E3: checkpoint intervention and short-horizon durability.}
E3 tests whether the candidate window identified in E1 is functionally consequential under
intervention. For a checkpoint $c_s$ at stage $s$, I inject a synthetic factual-association
signal by full fine-tuning the model for a fixed number of steps at a fixed learning rate. The
signal consists of fictional entity--value associations and is novel by construction. Held-out
probes use paraphrased prompts and are scored by multiple-choice accuracy and by the
log-probability margin between the correct value and the strongest distractor.

The experiment records three quantities. Uptake is the held-out signal score immediately after
injection relative to the checkpoint before injection. Clean retention is the score after a
fixed amount of ordinary continuation training on a small corpus that contains none of the
synthetic facts. Degradation-resistance is measured by fine-tuning the injected model on
conflicting associations under a sweep of poison budgets and recording how much of the signal
survives. Retention and degradation are reported both absolutely and normalized by uptake, so
that apparent durability is not confused with differences in how strongly each checkpoint
learned the signal during injection.

This design measures stage-dependent amenability to fine-tuning and short-horizon durability,
not lifetime persistence over the full remaining pre-training schedule. The learning rate,
injection budget, continuation budget, signal size, and held-out probe set are fixed across
stages. The first E3 run uses Pythia-160M and the factual-association signal only; larger
models and additional signal types are treated as follow-up checks.
